\section{The native SR bootstrap without an unshifted selector}
\label{sec:native-sr-bootstrap}

The ordinary SR proof can be carried out in the native Choice-DI
calculus before ordinary logical connectives have been reconstructed.
The missing ingredient is a restricted implication whose antecedent is
an equation between Boolean terms.  Such an antecedent permits successor
cancellation by a finite decoder; general successor cancellation is not
assumed.  This supplies precisely the contextual finite case analyses
used in the arithmetic proof.

\subsection{Primitive setting and the two immediate derived rules}

Use the stated native calculus with SR0/SRS, retaining its
Eq/$Q$/$D$ formulas, hypotheses, equality transport, paired and closure
elimination, and both case and literal-equation progress introductions.
Successor is the numerical abbreviation $Sa=R(0,a;u.0)$.
All rules in this section are syntactic derivability statements in that
native calculus; no ordinary disjunction elimination is presumed.

For any native formula $B(y)$ the native rules derive unconditional
numerical cases
\[
 \frac{\Gamma\vdash B(0)\qquad\Gamma\vdash B(Sz)}
      {\Gamma\vdash B(t)},\qquad z\text{ fresh for retained data}.
 \tag{NC}\label{eq:native-sr-NC}
\]
Indeed case introduction with an unused witness coordinate constructs
$D_{x,y}B(y)$ using the successor proof and ignoring its local base
assumption.  Paired elimination supplies $Q_yB(y)$, and closure
specialization with the given base proof supplies $B(t)$.
The analogous use of progress introduction gives contextual induction
on a literal equation, with the same retained-parameter restrictions.

Reflexivity requires no injectivity: the fixed computation
$R(1,0;u.t)=t$ can be used both as the major and the minor of transport
with motive $v=t$, yielding $t=t$.  Symmetry, transitivity, and numerical
congruence follow by the usual explicit equality motives.

Let $T(n,a,b)=R(n,a;u.b)$ and $Z(n)=T(n,0,0)$, as in the arithmetic
construction.  Under $0=1$, transport in the control coordinate gives
$T(0,0,s)=T(1,0,s)$, hence $1=s$, for every $s$.  Thus every equation
follows from the designated absurdity.  This atomic explosion is enough
for all arithmetic arguments below.  It does not use a selector with
unshifted endpoints.

\subsection{A finite successor decoder, before comparison}

The equations for $Z$ immediately give, by NC on $r$,
\[
                         Z(Z(Zr))=Zr.
 \tag{Z3}\label{eq:native-sr-Z3}
\]
Define $N(n)=Z(R(n,0;Z))$.  Direct computation and Z3 give
\[
             N(0)=0,\qquad N(Sn)=Z(N(n)),\qquad N(S^2n)=N(n).
\]
NC on $r$, using $N(0)=0$ and $N(1)=1$, also proves
\[
                              N(Zr)=Zr.
\]
Consequently the finite decoder
\[
                              p_2(w):=N(Sw)
\]
satisfies the unconditional identity
\[
                       p_2(S(Zr))=Zr.
 \tag{BD}\label{eq:native-sr-boolean-decoder}
\]
All these proofs are equation computations and unconditional NC.
Neither a Boolean-range disjunction nor contextual Boolean case
elimination was used to obtain BD.

In what follows a \emph{Boolean term} means syntactically a term $Zr$,
or a fixed numeral $0,1$ identified with $Z1,Z0$.  Both $N(n)$ and the
eventual comparator $L(x,y)$ are Boolean terms in this strict sense.
For Boolean $b,c$, applying $p_2$ to $Sb=Sc$ yields $b=c$ by BD.
Conversely congruence gives $Sb=Sc$ from $b=c$.
This is a local finite decoder, not general successor injectivity.

\subsection{Restricted implication and genuine contextual Boolean cases}

For Boolean $b,c$ and an arbitrary equation $E$ written $s=t$, define
the following native formula, with fresh binders and unused $x$:
\[
 \mathcal I(b=c,E)
 :=D_{x,y}\bigl(T(y,b,s)=T(y,c,t)\bigr).
 \tag{BI}\label{eq:native-sr-boolean-implication}
\]
Its zero matrix is $Sb=Sc$, equivalent to $b=c$ by BD; its successor
matrix is exactly $E$ by computation.  The native rules therefore derive
the following two rules:
\[
 \frac{\Gamma,b=c\vdash E}{\Gamma\vdash\mathcal I(b=c,E)},
 \qquad
 \frac{\Gamma\vdash\mathcal I(b=c,E)\qquad\Gamma\vdash b=c}
      {\Gamma\vdash E}.
 \tag{BII,BIE}\label{eq:native-sr-boolean-implication-rules}
\]
For introduction, choose an arbitrary unused witness, for instance zero.
The retained zero matrix $Sb=Sc$ supplies $b=c$ after applying $p_2$.
Insert the given premise derivation and rewrite its conclusion into the
successor matrix before closing case introduction.  For elimination,
paired elimination supplies the closure of that matrix.  Congruence
turns the given $b=c$ into its zero instance; specialization at $1$
then gives $E$.  The opened witness is unused and fresh.  Ordinary
assumption cut here means finite substitution of derivations, admissible
by induction on native proof trees.

Now let $b=Zr$ be a fixed Boolean term.  For an arbitrary retained native
context $\Gamma$ and an equation $E$, the following contextual case rule
is derivable:
\[
 \frac{\Gamma,b=0\vdash E\qquad\Gamma,b=1\vdash E}
      {\Gamma\vdash E}.
 \tag{BC}\label{eq:native-sr-boolean-cases}
\]
Choose a new numerical variable $w$ absent from $\Gamma,b,E,r$, and
apply unconditional NC to the native formula
\[
                       B(w):=\mathcal I(b=Zw,E).
\]
Its zero proof is BII applied to the supplied $b=1$ branch, since
$Z0=1$.  Its successor proof is BII applied to the supplied $b=0$
branch, since $Z(Sz)=0$.  NC gives $B(r)$.  Its antecedent $b=Zr$
is reflexive, so BIE gives $E$.

The retained context does not change during this argument.  In
particular it may contain a literal induction hypothesis mentioning
the original counter and the Boolean term $b$.  The varying variable
$w$ is fresh, distinct from that original counter.  This is the precise
reason BC can replace the ordinary proof's Boolean case analyses;
simply applying NC to a context that depends on its case variable
would have been invalid.

The consequent of $\mathcal I$ need only be an equation.  Although
$B(w)$ is compound, unconditional NC is available for every native
formula.  No selector for arbitrary formulas, unrestricted implication,
or trit case-elimination rule is required.

\subsection{The modulo-three cycle without a range disjunction}

Use the same numerical definitions
\[
 j(u)=T\bigl(u,0,T(Z(N(u)),1,0)\bigr),\qquad
 M(n)=j(j(R(n,0;j))).
\]
The fixed computations $j(0)=1$, $j(1)=2$, $j(2)=0$ use only the
already established initial values of $N$.
Instead of importing a logical statement that the range of $j$ is
three-valued, prove the single equation
\[
                              j^3(j(u))=j(u)
 \tag{J3}\label{eq:native-sr-j3}
\]
by unconditional NC on $u$.  At zero it is a fixed cycle computation.
At $u=Sz$, invoke BC on the Boolean term $N(Sz)$.  When its value is
zero, the inner selector gives $j(Sz)=0$; when its value is one, it
gives $j(Sz)=2$.  In both cases J3 is a fixed cycle computation.
The goals of this Boolean case use are equations.

Direct unfolding proves $M(0)=0$ and $M(Sn)=j(M(n))$.
Since $M(n)$ is itself of the form $j(v)$, J3 gives
\[
                              M(S^3n)=M(n).
\]
This obtains the exact modulo-three equation used by the remaining
construction without assuming a native three-way disjunction or a
successor decoder on arbitrary numbers.

\subsection{Transferring the arithmetic proof, step by step}

Every subsequent step of the comparator and predecessor proofs now has
a native derivation with the same numerical terms.
The dependencies are as follows.

\begin{enumerate}
\item The doubling and tripling residue invariants are literal equations.
Their successor steps rewrite the update directly with their respective
induction equations $N(d_2(n))=0$ and $M(d_3(n))=0$.  They use the
two- and three-period equations just proved, and need no range cases.

\item The ramp invariant $M(G(n))=N(n)$ and the two scaling identities
use Boolean case splits only on $N(n)$.  Replace each of those ordinary
splits by BC; retain the original literal induction equation in $\Gamma$.
The finite update values at residues zero and one are direct computations.
The unused residue-two reset is also a direct fixed computation.
No trit elimination rule is needed.

\item The factor-map identities split only $N(u)$.  Seed rotation,
iterator factorization, and the zero row of comparison are equation
inductions with fixed passive parameters.  The comparator recurrences
are therefore native equation theorems.

\item Comparison persistence is a contextual equational derivation from
the premise $L(x,y)=1$.  Split the Boolean $N(R(y,E(x);H))$ by BC.
Its value one contradicts that retained premise by fixed computation;
atomic explosion handles this branch.  Its value zero makes $H$ fix the
state and proves the next comparison equation.  Other conditional
comparison facts use the same pattern, with Boolean $L$ terms.

\item The bounded-cycle run is obtained by induction on the same literal
gate equation
\[
 T(L(t,m),0,J_m(t))=T(L(t,m),0,t).
\]
Its successor proof uses BC on $L(St,m)$ while retaining that literal
induction equation.  The downward comparison fact is used as its
already established contextual equation derivation.  The period and
doubled-trajectory arguments are then pure equation inductions.
Consequently the native calculus proves $P(0)=0$ and $P(Sa)=a$.
\end{enumerate}

This is a finite proof transformation: replace each recorded Boolean
case with the displayed BC derivation, replace any impossible Boolean
branch with native atomic explosion, and use the native derived EI
for each displayed literal induction.  All remaining steps are
computation, equality transport, substitution, or unconditional NC.
No part of the proof asks for induction on a native compound formula.

\subsection{Closing the full logical and arithmetic equivalence}

With the native predecessor theorem available, the finite macro
\[
 I(n,a;u.h):=P\bigl(R(n,a;v.S(h[Pv/u]))\bigr)
\]
has the ordinary initialized-iterator equations by native computation
and transport.  In particular $C(n,a,b)=I(n,a;u.b)$ is now an
unshifted selector.  At this point the established Choice-DI gate,
padding, formula-selector, and logical-macro reconstruction applies
without circularity.  It gives the complete intuitionistic first-order
logical interface, coverage and separation, and the native formula
return translations.  The ordinary SR arithmetic and composite-uniqueness
proofs then apply to this reconstructed interface.

Conversely every native SR rule is sound in the ordinary SR theory
under the usual readings of $Q$ and $D$: case introduction uses numerical
coverage, and progress introduction uses only literal-equation induction.
The domains and numerical equality are unchanged, and the familiar
$Q$/$D$ return equivalences are available after $C$ has been derived.
Thus the direct native SR calculus itself is equiconsistent and
bi-interpretable with full first-order iPRA.  A serialized successor
node may be retained as the charged abbreviation $R(0,a;u.0)$; its
presence is an implementation convenience, not a new numerical axiom.
